% English manuscript, embedded-inclusion version (October 2026).
% Values in guillemets are filled from the records by fill_embedded.py.
\documentclass[11pt,a4paper]{article}
\usepackage[T1]{fontenc}
\usepackage[utf8]{inputenc}
\usepackage[english]{babel}
\usepackage[margin=25mm]{geometry}
\usepackage{amsmath,amssymb,graphicx}
\usepackage{placeins}
\usepackage[numbers,sort&compress]{natbib}
\usepackage[hyphens]{url}
\usepackage[hidelinks]{hyperref}
\ifx\pdfmapfile\undefined\else
\pdfmapfile{+cm-super-t1.map}\pdfmapfile{+cm-super-ts1.map}\fi
\setlength{\emergencystretch}{2em}
\renewcommand{\topfraction}{0.85}
\renewcommand{\textfraction}{0.10}
\renewcommand{\floatpagefraction}{0.75}

\title{Stresses in aluminium around an Al$_{13}$Fe$_4$ inclusion with a given
magnetostrictive strain}
\author{I.~Mikhailovskiy\thanks{Corresponding author.
\texttt{ilyamihailovsy@gmail.com}} \and D.~Pshonkin}
\date{Moscow Polytechnic University, Moscow, Russia}

\begin{document}
\maketitle

\begin{abstract}
Molecular dynamics at 0~K gives stresses around a spherical
Al$_{13}$Fe$_4$ inclusion of radius 2.5~nm inside aluminium, with a given
magnetostrictive strain along the magnetic field. Near it the matrix is
compressed by «POLE_PEAK»~MPa along the field and stretched by «EQ_PEAK»~MPa across
it; shear stresses reach «SHEAR_PEAK»~MPa. The response is elastic: even
tenfold strain forms no dislocations, and a dislocation above the
inclusion shifts «DISL_SHIFT»~\AA{} and returns on unloading. Compliant,
distorted aluminium at the interface makes the atomistic field about
«A_PCT»\% of the elastic solution. If this loss vanishes for micron
particles (untested), this strain gives elastic stresses above the alloy
yield stress in a shell about one radius thick. With the field on, the
observed creep increase would need a third of the measured Al--Mg--Si
activation volume, or 6--7 times the measured Fe--Al magnetostriction.
Why the effect outlasts the field remains unexplained.
\end{abstract}

\noindent\textbf{Keywords}\quad molecular dynamics; magnetostriction;
inclusions; internal stress; aluminium alloys

\section{Introduction}

Al--Mg--Si alloys containing iron-rich inclusions creep faster after they
have been held in a static magnetic field (creep is slow deformation under a
constant load), and the effect persists after the field is removed
\cite{Skvortsov2019JAP,Skvortsov2019PSS,Friha2024JMMM,Nikolaev2025SciRep}.
One proposed explanation is magnetostriction of the inclusions, the change of
their shape in the field. The matrix resists the change, and stresses appear
around the inclusion. These stresses can move dislocations, the line defects
whose motion carries plastic deformation. Ref.~\cite{Friha2024JMMM} estimates the
stress at the inclusion boundary as 147~MPa in a field of 0.7~T. This estimate
does not say how the stress is distributed around the
particle, how fast it falls with distance, or how much of it acts on
dislocations.

Here the stress field around one Al$_{13}$Fe$_4$ inclusion in aluminium is
calculated. The strain of the inclusion is prescribed: molecular dynamics
follows the atoms with an interatomic potential, a formula for the energy of
interacting atoms, and this potential does not describe magnetism, so the
magnetostriction cannot come out of the calculation itself. Its magnitude
corresponds to the 147~MPa estimate \cite{Friha2024JMMM}. The atomistic
result is compared with the elastic solution, and the field is then used to
estimate the force on dislocations and the increase in creep rate.

\section{Model and method}

\subsection{The cell}

The cell is a periodic box of aluminium, $12.6\times12.4\times12.6$~nm.
A spherical Al$_{13}$Fe$_4$ inclusion of radius $R=2.5$~nm sits at its centre
(Fig.~\ref{fig:cell}a). The inclusion is cut from an Al$_{13}$Fe$_4$ crystal
with the lattice of the crystal-structure data \cite{Liu2024CIF}; aluminium
atoms closer than 2.2~\AA{} to an inclusion atom are removed. The cell holds
118\,861 atoms, 114\,374 in the matrix and 4\,487 in the inclusion. Aluminium surrounds the inclusion on all sides. The boundary
conditions are periodic along all three axes, so the cell represents a
regular array of inclusions 12.6~nm apart along $x$ and $z$ and 12.4~nm along $y$; their volume fraction is 3.3\%.

A sphere was chosen as the simplest compact shape; the actual shape of the
inclusions is not reproduced. The stress
around a sphere falls with the distance $r$ from its centre as
$r^{-3}$ and is small two or three radii away, so a cell of $10^5$ atoms is
large enough. A sphere also has an exact elastic solution
\cite{Eshelby1957,Mura1987} against which the atomistic result can be checked.

The cell axes are $x\parallel[1\bar10]$, $y\parallel[11\bar2]$ and
$z\parallel[111]$. The $(111)$ planes on which dislocations glide in
aluminium are then horizontal, and the slip direction $[1\bar10]$ lies along
$x$. Atoms interact through the second-nearest-neighbour modified embedded-atom
potential for Al--Si--Mg--Cu--Fe \cite{Jelinek2012}; the calculations use
LAMMPS \cite{Thompson2022LAMMPS}. The elastic constants of aluminium with
this potential, $C_{11}=110.5$, $C_{12}=60.9$ and $C_{44}=28.4$~GPa, are close
to the measured ones.

\subsection{The given magnetostrictive strain}
\label{sec:strain}

The magnetic field points along $x$. Magnetostriction lengthens the inclusion
along the field and shortens it across, without changing its volume
(Fig.~\ref{fig:cell}b):
\begin{equation}
\varepsilon_{xx}=\varepsilon,\qquad
\varepsilon_{yy}=\varepsilon_{zz}=-\varepsilon/2,\qquad
\varepsilon=1.94\times10^{-3}.
\label{eq:strain}
\end{equation}
This $\varepsilon$ is the 147~MPa estimate of Ref.~\cite{Friha2024JMMM}
divided by a Young's modulus of 75.7~GPa; the modulus of aluminium along
$[111]$ with the potential used here is 75.9~GPa. Ref.~\cite{Friha2024JMMM}
used 70~GPa, with which 147~MPa corresponds to a strain of
$2.1\times10^{-3}$. The magnetostriction of Al$_{13}$Fe$_4$
inclusions has not been measured, so $\varepsilon$ is an input of the
calculation. Every stress below is proportional to $\varepsilon$; for another
strain $\varepsilon'$ it is multiplied by $\varepsilon'/0.194\%$ (by 1.08 for
$2.1\times10^{-3}$).

Each inclusion atom is moved by
$\mathbf u=\boldsymbol\varepsilon(\mathbf r-\mathbf r_c)$, where $\mathbf r$ is
the position of the atom before the move, $\mathbf r_c$ the inclusion centre
and $\boldsymbol\varepsilon$ the strain (\ref{eq:strain}). The inclusion atoms
then stay fixed, while the aluminium atoms move until the forces on them vanish
(energy minimisation by conjugate gradients). This gives the inclusion the
shape it is assumed to take inside the matrix under the field. The arrows in
Fig.~\ref{fig:cell}b show the forces with which the inclusion acts on the
aluminium, taken from the radial stress at the interface in the elastic
solution for a single sphere: outward at
the ends along the field and inward on top, at the bottom and at the sides.
In a real inclusion magnetic forces would maintain this shape, if a
magnetostriction of this size exists.

The strain (\ref{eq:strain}) is the strain of the inclusion inside the matrix.
A free particle would lengthen more in the same field, because no matrix
would hold it back. If 0.194\% were the free magnetostriction of an elastic
particle, every matrix stress below would be multiplied by
$A=\beta k/(1+\beta(k-1))$, where $k$ is the ratio of the shear moduli of
inclusion and aluminium and $\beta=0.46$ is a coefficient of the elastic
solution for a sphere that depends only on Poisson's ratio (Supplementary S3);
for an inclusion
two to two and a half times stiffer than aluminium, $A=0.63$--$0.68$.

The reference cell, without strain, is prepared once. The inclusion atoms are
fixed from the start: with this potential the equilibrium lattice of
Al$_{13}$Fe$_4$ is 8\% larger in volume than the measured one, and a free
inclusion would swell. The matrix is held at 300~K for 10~ps, cooled to
0.5~K over 10~ps and relaxed at 0~K. The cell size is then adjusted once so
that the mean stress in the aluminium far from the inclusion vanishes. The
cells with strain $+\varepsilon$ and $-\varepsilon$ are built from the
reference cell by moving the inclusion atoms; only the matrix relaxes in
them, at fixed cell size. Half the sum of the responses to $+\varepsilon$ and
$-\varepsilon$ shows how far the response departs from linear
(Supplementary S2).

\FloatBarrier
\begin{figure}[!htbp]
\centering
\includegraphics[width=\linewidth]{fig_cell}
\caption{The model. (a)~Section of the cell by the $xz$ plane through the
inclusion centre, a 1~nm thick slab (OVITO \cite{Stukowski2010OVITO}); grey dots are aluminium atoms of the
matrix, orange and red ones Al and Fe atoms of the inclusion. (b)~The given
strain of the inclusion (exaggerated): dashed circle, original shape; filled,
shape with the strain (\ref{eq:strain}). Arrows: forces with which the
inclusion acts on the aluminium in the single-sphere elastic solution,
outward at the ends along the field and inward on top and at the bottom;
arrow length is proportional to force per unit surface area, and the numbers
give this force along and across the field.}
\label{fig:cell}
\end{figure}
\FloatBarrier

\subsection{How the stresses are obtained}

The stress at each matrix atom is its virial divided by the atomic volume of
aluminium, 16.61~\AA$^3$. The virial of an atom is the sum of the products of
the forces from its neighbours and the vectors joining it to them; in a uniform
crystal the virials summed and divided by the volume give the ordinary stress. The
quantity used is the difference between the strained cell and the reference
cell, that is, the stress added by the magnetostriction. Tension is
positive and compression negative. Three quantities are used. The normal
stress along the radius, $\sigma_{rr}$, shows how hard the matrix pushes on
the inclusion or pulls it; the 147~MPa estimate is of this kind. The
equivalent (von Mises) stress, $\sigma_{\text{eq}}$, combines all components
except the hydrostatic pressure into one number that is compared with the
yield stress, the stress at which the alloy starts to deform plastically. The shear stress
$\tau=\sigma_{xz}$ acts on the horizontal slip plane along $x$; it is the
stress that pushes an edge dislocation with Burgers vector along $x$. The
Burgers vector is the shift of the crystal left behind when the dislocation
passes; for an edge dislocation it is perpendicular to the line.

Values are averaged over atoms in spherical shells 2~\AA{} thick inside narrow
cones about chosen directions: along the field (within 10$^\circ$), across it
(within 10$^\circ$ of the plane $x=0$), and at 25$^\circ$, 45$^\circ$ and
67$^\circ$ to the field (bands of $\pm7.5^\circ$). The distance $d=r-R$ is
measured from the inclusion surface. Apart from the weak elastic anisotropy
of aluminium (the dependence of its elasticity on direction) the field is symmetric about the $x$ axis, so for the maps in
Fig.~\ref{fig:map} the stresses are averaged over rotations about that axis.

\subsection{The elastic solution}

The atomistic result is compared with two elastic solutions. The first is the
exact solution for a single sphere in an infinite isotropic matrix with the
same inclusion strain \cite{Eshelby1957,Mura1987,JuSun1999}. The second is
obtained for the same periodic cell by a Fourier method
\cite{MoulinecSuquet1998} with the cubic elastic constants of aluminium from
the potential. It includes the neighbouring inclusions of the periodic array
and the fixed cell size, so it can be compared with the atomistic result
directly. The first shows what the field of a single inclusion in a large
grain would be.

\section{Results}

\subsection{How the stress is distributed}

Fig.~\ref{fig:map} shows how the added stress is distributed around the
inclusion. Along the field, at the ends of the inclusion, the matrix is
compressed: the inclusion pushes it apart. Across the field the matrix is in
tension: the inclusion draws back from it. The shear stress on the slip planes
vanishes on the $x$ and $z$ axes and is largest in four inclined lobes, at
about 25$^\circ$ and 67$^\circ$ to the field, changing sign from lobe to lobe.
The atomistic map and the elastic map for the same cell agree in shape, but
the atomistic stresses are about one fifth smaller
(Section~\ref{sec:decay}). Within «LAYER»~\AA{} of the interface the shape
differs as well, because a continuum does not describe the atomic structure
of the boundary.

\FloatBarrier
\begin{figure}[!htbp]
\centering
\includegraphics[width=\linewidth]{fig_map}
\caption{Stresses added by the strain of the inclusion in the $xz$ plane
through its centre: equivalent (von Mises) stress (top) and shear stress on
the slip plane, $\tau=\sigma_{xz}$ (bottom). Left, atomistic calculation;
right, elastic solution for the same periodic cell. The field points along
$x$; the circle is the inclusion surface, and the white ring around it
(3~\AA{} of matrix) is left blank. Dotted lines in the lower left panel:
directions at 25$^\circ$ and 67$^\circ$ to the field.}
\label{fig:map}
\end{figure}
\FloatBarrier

\subsection{How the stress falls with distance}
\label{sec:decay}

Fig.~\ref{fig:decay} gives the profiles along the chosen directions. The
atomistic values quoted here are for shells 5--10~\AA{} from the surface;
closer in, the narrow cones along the field hold only 15--20 atoms per shell,
and the values scatter by up to 25~MPa. Along the field the normal stress is
compressive, about «MD_POLE_PEAK»~MPa, which is close to the
147~MPa estimate of Ref.~\cite{Friha2024JMMM}. The agreement is expected,
because $\varepsilon$ was derived from that estimate. Across the field the
normal stress is tensile, about «MD_EQ_PEAK»~MPa. At the sides of the
inclusion, near the plane $x=0$, the stress along the field is also large:
as the inclusion lengthens, it pulls the aluminium next to it along, and
$\sigma_{xx}$ right at the surface there reaches «MD_EQ_SXX»~MPa. The shear
stress on the slip planes reaches «MD_SHEAR_PEAK»~MPa in the lobes.

One radius from the surface ($r=2R$) the normal stress along the field falls
to about 54~MPa and the shear stress in the lobes to about 19~MPa; two radii away ($r=3R$) they are about
18 and 7~MPa (single-sphere solution).

From $d\approx$~«FIT_D»~\AA{} on, the atomistic stresses are close to the
elastic solution for the same cell times «A_FIT»; after this scaling they
differ from it by «MED_PCT»\% on average. The factor depends a little on the
direction: «A_POLE» along the field, «A_EQ» across it, and «A_L25» and «A_L67»
for the shear stress in the 25$^\circ$ and 67$^\circ$ lobes. The atomic
displacements, which do not depend on how the stress is computed, are reduced
in the same way, and more than 8~\AA{} from the interface the stress of each
atom equals the elastic constants of aluminium times its strain to within 3\%.
The matrix therefore deforms less than in the elastic solution with a
perfectly bonded interface. The deficit builds up within 5--8~\AA{} of the
interface: in the first aluminium layer the displacement is «KIN_RAD»\% of the
elastic one along the surface normal and «KIN_TAN»\% along the surface, at
4--6~\AA{} it is already «D46R»\% and «D46T»\%, and further out it changes
little. In this layer the aluminium lattice is distorted, the first layer has
half the bulk density, and the atoms carry less stress for the same strain
(Supplementary S2).

If the compliance of this layer does not depend on particle size, the
weakening scales inversely with the radius and would be about «LOSS_UM»\% for
a particle of radius 2.5~$\mu$m, so the elastic solution is used below for
real particles. This assumption is untested: one inclusion radius and one
interface (aluminium atoms within 2.2~\AA{} of the inclusion removed) were
computed, at 0~K. Sliding at the interface at service temperature would give
a weakening that does not depend on particle size.

The solution for the cell differs from the single-sphere solution: it has
the cubic elasticity of aluminium, a fixed cell size and neighbouring
inclusions. The shear stress in the lobes differs between them by at most
2~MPa for $3\le d\le20$~\AA, while the normal stresses in the cell are
13--17\% larger already at the surface and 30--44\% larger at $r=2R$. The volume
fraction of inclusions in the alloy, 0.25\%, is more than ten times smaller
than in the cell, so the decay far from the inclusion is shown by the
single-sphere solution.

The field of a single sphere depends only on $r/R$, so the elastic solution carries over
to inclusions of real size: for a particle 5~$\mu$m across ($R=2.5~\mu$m) the
distances in Fig.~\ref{fig:decay} are multiplied by a thousand (upper scale).
At the given strain the equivalent stress exceeds the yield stress of the alloy, 40--60~MPa
\cite{Friha2024JMMM}, out to $r\approx\text{«SHELL_VM»}R$. The largest shear stress on the
twelve slip systems (a slip system is a slip plane together with a slip
direction in it) exceeds the critical value for slip, 13--20~MPa, out to
$r\approx\text{«SHELL_SLIP»}R$. This critical value is the yield stress divided by the
Taylor factor 3.06 \cite{StollerZinkle2000}, which accounts for the different
orientations of the grains in a polycrystal. For a 5~$\mu$m particle the shell where the equivalent stress exceeds the yield stress is «THICK_VM_UM»~$\mu$m
thick. Inside this shell the elastic solution overestimates the stresses: in the
alloy, dislocations already present would start to move, the matrix would
flow plastically and part of the stress would relax. The atomistic
calculation does not test this: its matrix is pure aluminium without
pre-existing dislocations, and at 0~K no new ones form in it even at a lobe
shear stress of «X10_SHEAR»~MPa (strain $10\varepsilon$, Supplementary S2).

\FloatBarrier
\begin{figure}[!htbp]
\centering
\includegraphics[width=\linewidth]{fig_decay}
\caption{Added stresses against the distance $d$ from the inclusion surface:
(a)~along the field, (b)~across the field, (c)~shear stress in the
25$^\circ$ and 67$^\circ$ lobes. Symbols, atomistic calculation (mean over a
2~\AA{} shell and its error); solid lines, elastic solution for the same cell;
dashed lines, single sphere; hatched, the layer at the interface where they
differ in shape. Away from the interface the symbols lie closer to zero than the solid lines:
the atomistic field is about «A_PCT»\% of the same-cell solution. The
horizontal band is the yield stress of the
alloy, 40--60~MPa. Upper scale: distance for a particle of radius 2.5~$\mu$m.
(d)~Magnitude of the single-sphere stress far from it against $r/R$ on logarithmic scales;
dotted, the slope $r^{-3}$.}
\label{fig:decay}
\end{figure}
\FloatBarrier

\subsection{The force on dislocations}

A shear stress $\tau$ exerts on a dislocation with Burgers vector
$b=2.86$~\AA{} a force $\tau b$ per unit length, directed along the slip plane
(the Peach--Koehler force \cite{HirthLothe1982}). The dislocation moves if this
force exceeds what holds it back: the lattice resistance (the Peierls stress,
the least stress at which a straight dislocation in a pure crystal at 0~K
moves to the next position), solute atoms and other dislocations.

Fig.~\ref{fig:glide} shows the shear stress of the inclusion field averaged
along a straight line on which an edge dislocation can lie. The line runs
parallel to $y$ at a height $h$ above the inclusion surface. The dislocation is as long as the cell period along $y$,
12.4~nm. The stress changes sign when $x$ is replaced by $-x$: on one side of
the inclusion it pushes the dislocation forward, on the other side back. The averaged stress is largest,
«LINE_PEAK»~MPa, at $h=$~«LINE_H»~\AA. On the short part of the line right
above the inclusion the stress is larger, up to «LOCAL_PEAK»~MPa, and this part
could bow out even when the average is too small to move the whole line. The
line average depends on the spacing of the inclusions in the model
(12.4~nm along the line) and does not carry over directly to the alloy.

In an Al--Mg--Si alloy dislocations are held by Mg and Si atoms and by
precipitates, fine particles of Mg--Si compounds inside the grains; the critical shear stress of the alloy, from its yield stress,
is 13--20~MPa. In the single-sphere solution at the given strain the field exceeds this
only in a shell «THICK_SLIP» radii thick next to the inclusion. Farther out the field is below this
value and cannot by itself push dislocations past their obstacles; it can
only speed up or slow down their thermally activated motion, in which thermal
vibrations of the atoms carry them past the obstacles (see the creep
estimate below).

«DISL_TEXT»

\FloatBarrier
\begin{figure}[!htbp]
\centering
\includegraphics[width=0.75\linewidth]{fig_glide}
\caption{Shear stress averaged along a straight edge dislocation 12.4~nm long
that runs above the inclusion at a height $h$ above its surface, against the
position $x$ of the dislocation. Symbols, atomistic calculation; lines,
elastic solution for the same cell; shaded, the extent of the inclusion
along $x$. Right scale: force per unit length of the dislocation, $\tau b$.}
\label{fig:glide}
\end{figure}
\FloatBarrier

\subsection{An estimate for creep}

The link between the stress field and the creep rate is estimated with a model
of thermally activated glide \cite{KocksArgonAshby1975,CaillardMartin2003}.
The estimate applies to creep while the field is on and the stresses around
the inclusions exist; the observed increase was measured after the field was
removed (see the Discussion).
A dislocation passes its obstacles with the help of thermal fluctuations. A
shear stress $\tau$ makes forward jumps $\exp(V^*\tau/k_{\mathrm B}T)$ times
more likely and backward jumps less likely by the same factor; $V^*$ is the
activation volume, a measure of how strongly the barrier depends on stress,
and $k_{\mathrm B}T$ is the thermal energy. If dislocations around an
inclusion meet stresses of both signs equally often, their mean velocity grows
by the factor $\cosh(V^*\tau/k_{\mathrm B}T)$, the mean of the factors
$\exp(V^*\tau/k_{\mathrm B}T)$ and $\exp(-V^*\tau/k_{\mathrm B}T)$ for the two
signs. The excess of this factor over one, integrated over the matrix around
one inclusion and multiplied by the number of inclusions per unit volume,
$f/V_p$, gives the relative increase in creep rate (derived in
Supplementary S5),
\begin{equation}
H=\frac{f}{V_p}\int_{r>R}
\Bigl\langle\cosh\frac{V^*\tau_s(\mathbf r)}{k_{\mathrm B}T}-1\Bigr\rangle
\,\mathrm dV ,
\label{eq:creep}
\end{equation}
where $f$ is the volume fraction of inclusions, $V_p$ the volume of one
inclusion, $\tau_s$ the shear stress on slip system $s$ (a slip plane together
with a slip direction in it), and the angle
brackets an average over the twelve slip systems and over grain orientations.
The elastic field of a single sphere goes into the integral (the atomistic
weakening is taken as negligible for micron-size particles,
Section~\ref{sec:decay}); it depends only
on $r/R$, so $H$ does not depend on particle size. The values used are
$T=300$~K and $f=0.246\%$; the latter corresponds to 0.35~wt\% of inclusions
\cite{Friha2024JMMM}.

In Fig.~\ref{fig:creep} the solid curve shows the amplitude $\tau_m$ needed for
the 25\% increase in creep rate observed after holding in the field
\cite{Skvortsov2019JAP} at different $V^*$ if the shear stress
falls from the inclusion surface as $\tau_m(R/r)^3$; this amplitude falls as
the activation volume grows. The dashed curve shows the amplitude $\tau_m$ for
which such a field gives the same increase as the computed field
(Supplementary S5). The curves cross at $V^*\approx«VSTAR_GRAINS»\,b^3$. At $V^*=70\,b^3$, measured for
an Al--Mg--Si alloy \cite{Soula2022JALCOM}, the computed field would raise the
rate $10^8$ times; the model is then outside its range, and a much smaller
strain would be enough for the observed 25\%. Because $H$ depends only on the
product $V^*\varepsilon$, at $V^*=70\,b^3$ an inclusion strain of about
«EPS_AT_70»\% is needed. The measured magnetostriction of Fe--Al alloys does not
exceed 0.01\% \cite{Hall1959,BormioNunes2012}, six to seven times less.

\FloatBarrier
\begin{figure}[!htbp]
\centering
\includegraphics[width=0.7\linewidth]{fig_creep}
\caption{Shear-stress amplitude needed for a 25\% increase in creep rate
(solid line) and the amplitude equivalent to the elastic single-sphere field of the
inclusion with a strain of 0.194\% (dashed), against the activation volume
$V^*$; $T=300$~K, $f=0.246\%$. Dotted line: $V^*=70\,b^3$ for Al--Mg--Si
\cite{Soula2022JALCOM}.}
\label{fig:creep}
\end{figure}
\FloatBarrier

\section{Discussion}

The calculation gives the stresses that an inclusion of a given shape
produces in aluminium. Whether the inclusions take this shape in a field is
not known: the magnetostriction of Al$_{13}$Fe$_4$ inclusions has not been
measured, and strong ferromagnetism has not been established for bulk
Al$_{13}$Fe$_4$ \cite{Albedah2015Mossbauer,Avers2025Crystals}. If the real
strain of the inclusions is close to the magnetostriction of Fe--Al alloys,
the stresses are twenty times smaller than computed, and the shell where they
exceed the yield stress disappears. From the single-sphere solution, the largest shear
stress on the slip systems at the particle surface exceeds the critical value of the alloy once the
inclusion strain is above «EPS_SLIP_LO»--«EPS_SLIP_HI»\%, and the equivalent
stress exceeds the yield stress above «EPS_VM_LO»--«EPS_VM_HI»\%. These are
strains of the inclusion inside the matrix; the free magnetostriction of an
elastic particle would have to be $1/A$ times larger
(Section~\ref{sec:strain}).

The inclusion is held in its given shape, so its own elasticity does not
enter the problem. The potential describes the
elasticity of aluminium well but overestimates the lattice of Al$_{13}$Fe$_4$,
which is why the inclusion is held at the measured lattice. The calculation is
static (energy minimisation without thermal motion of the atoms) and done at
0~K for one inclusion of one size (radius 2.5~nm) in a regular array;
the elastic solution for the same cell includes the field of the neighbouring
inclusions. The matrix is pure aluminium without Mg and Si atoms,
precipitates or pre-existing dislocations, so the yield stress of the alloy
is taken from the literature \cite{Friha2024JMMM} and plastic flow near the
inclusion is not modelled.

The main limitation concerns the persistence of the effect. In the experiments creep is
measured after the field is removed, whereas the elastic stress of
magnetostriction disappears with the field. If the field moves dislocations
near the inclusions, the displaced dislocations may stay in their new positions
after the field is removed; the computed shell where the stress exceeds the
yield stress shows where to look for such changes.
In the one case computed, a straight dislocation at 0~K, the dislocation
returned when the strain was removed; the calculation does not show whether
such a rearrangement of defects persists in the alloy or how it affects
subsequent creep.

\section{Conclusions}

An Al$_{13}$Fe$_4$ inclusion of radius 2.5~nm given a magnetostrictive strain
of 0.194\% along the field compresses the aluminium matrix along the field by
about «MD_POLE_PEAK»~MPa and stretches it across the field by about
«MD_EQ_PEAK»~MPa. The shear stress on the slip planes reaches
«MD_SHEAR_PEAK»~MPa in inclined lobes. In the static atomistic calculation at 0~K the matrix responds
elastically: no dislocations form even at ten times the strain, and a dislocation above the
inclusion is moved by the field by «DISL_SHIFT»~\AA{} and returns when the
strain is removed. The atomistic field is about «A_PCT»\% of the elastic
solution for a perfectly bonded inclusion; the weakening comes from a
compliant layer of distorted aluminium at the interface. If its compliance
does not depend on particle size, the weakening is negligible for
micron-size particles, but one inclusion radius at 0~K does not test this. Far from the
inclusion the field falls as $(R/r)^3$; in the elastic solution at the
given strain the matrix is
stressed beyond the yield stress of the alloy in a shell about one radius
thick, «THICK_VM_UM»~$\mu$m for a 5~$\mu$m particle. In this shell the alloy could
deform plastically, that is, by irreversible motion of dislocations already
present; this case was not computed atomistically. In a model of thermally
activated glide this field, if it acted during creep, would give the observed
increase in creep rate for an
activation volume of about «VSTAR_GRAINS»$\,b^3$; at the measured $70\,b^3$ a strain of «EPS_AT_70»\% would
suffice, six to seven times the measured magnetostriction of Fe--Al alloys.
The calculation does not explain why the effect persists after the field is
removed.

\section*{Author contributions}

I.~Mikhailovskiy developed the method and software, ran the calculations,
analysed the results, made the figures and wrote the text. D.~Pshonkin
conceived the study, supervised it, and reviewed and edited the text.

\section*{Data and code availability}

The LAMMPS inputs, the programs that build the cell, analyse the results and
draw the figures, the reference and strained configurations with atomic
stresses, and the result records are available at
\url{https://github.com/iluuua/magnetostriction-plasticity-bounds}.

\section*{Conflict of interest}

The authors declare no conflict of interest.

\bibliographystyle{unsrtnat}
\bibliography{references}
\end{document}
